\scanpage{25}

les $Dt_i=0$. \deleted{\unsure{Donc|On}} Prenons une \uline{base} de $\Omega^1_{K/k}$ sur $K$, \unsure{ou soit|on en} $(\omega_i)_{i\in I}$ donc $\Omega^1_{K/k}\otimes_K A\simeq A^{(I)}\subset A^I$. En tant que dérivation dans $A^I$, $D$ est donnée par ses composantes $D_i$, qui sont uniquement déterminées

\[ D_{\unsure{i|\vphantom{i}}}\left(\sum_\lambda a_\lambda t^\lambda\right)=\sum_\lambda\deleted{\unsure{D|d|D_i}}d_{K/k}(a_\lambda)t^\lambda \]

Il reste à prouver qu'on peut trouver des $f=\unsure{\sum|(}a_\lambda t^\lambda$ telles que les $(D_i(f))_{i\in I}\notin A^{(I)}$. \deleted{Il suffit} \inserted{Choisissons} \deleted{pour ceci} \inserted{pour ceci} \deleted{que les $D_i(a_{\unsure{n|\lambda}})$ pour tt} \deleted{de trouver} une \uline{suite infinie} $i_0,i_1,\ldots i_n$ d'indices distincts de $I$, \deleted{tels que} pour tt $n$ il existe $a_n\in K$ tel que \deleted{$D_{i_n}(a_n)\ne0$} \deleted{$D_{\unsure{i_n|n}}(a_n)=0$}, \deleted{Or ceci est possible} $d_{K/k}\deleted{\unsure{a^n|a_n}}a_n\ne0$, il suffit alors de prendre $f=\sum a_n t^n$. \ldots\ldots
% archive sheet number: 24
